\customchapter{ABSTRACT}

\begin{center}
\large \vspace{-1.5cm} \textbf{Mitigating evolutionary resistance in glioblastoma using adversarial multi-agent reinforcement learning}
\end{center}

Glioblastoma multiforme develops evolutionary resistance that limits the duration of tumor control achieved with temozolomide. This thesis proposes GBMARL, an adversarial multi-agent reinforcement learning framework that models therapy and tumor evolution as two agents interacting over a simulator of sensitive and resistant tumor populations. The therapy policy is trained with MAPPO under centralized training and decentralized execution (CTDE), and is compared with IPPO, no treatment, maximum tolerated dose, and the adaptive Gatenby heuristic.

The simulator was calibrated by integrating DepMap Public 26Q1 and GDSC2 release 8.5. The data join recovered 34 GBM cell lines with sensitivity assays and 24 lines with complete expression data. For temozolomide, the estimated sensitivity gap between the sensitive and resistant populations was approximately 9.1-fold. Evaluation uses a combined time-to-progression endpoint: an episode ends when tumor burden crosses its control threshold or when the resistant fraction becomes a majority, while recording the failure mode.

In the nominal calibrated comparison, fifteen MAPPO policies achieved a median of 34 days [27, 41] and exceeded Gatenby in 13/15 seeds; fifteen IPPO policies achieved a median of 31 days [23, 34] and exceeded Gatenby in 12/15 seeds. The paired MAPPO--IPPO difference was exploratory ($p=0.013759$). Sensitivity analysis identified a major fragility to reductions in carrying capacity $K$: at $K=0.8K_0$, the MAPPO median fell to 25 days. The results therefore provide in-silico evidence of adaptive behavior under the specified model, not evidence of clinical efficacy.

\noindent \textbf{Keywords:} glioblastoma; adaptive therapy; evolutionary resistance; multi-agent reinforcement learning; MAPPO-CTDE; temozolomide.
